\documentclass[12pt, titlepage]{article}

\usepackage{booktabs}
\usepackage{tabularx}
\usepackage{hyperref}
\hypersetup{
    colorlinks,
    citecolor=black,
    filecolor=black,
    linkcolor=red,
    urlcolor=blue
}
\usepackage[round]{natbib}

\input{../Comments}
\input{../Common}

\begin{document}

\title{Verification and Validation Report: \progname} 
\author{\authname}
\date{\today}
	
\maketitle

\pagenumbering{roman}

\section{Revision History}

\begin{tabularx}{\textwidth}{p{3cm}p{2cm}X}
\toprule {\bf Date} & {\bf Version} & {\bf Notes}\\
\midrule
Date 1 & 1.0 & Notes\\
Date 2 & 1.1 & Notes\\
\bottomrule
\end{tabularx}

~\newpage

\section{Symbols, Abbreviations and Acronyms}

\renewcommand{\arraystretch}{1.2}
\begin{tabular}{l l} 
  \toprule		
  \textbf{symbol} & \textbf{description}\\
  \midrule 
  T & Test\\
  \bottomrule
\end{tabular}\\

\wss{symbols, abbreviations or acronyms -- you can reference the SRS tables if needed}

\newpage

\tableofcontents

\listoftables %if appropriate

\listoffigures %if appropriate

\newpage

\pagenumbering{arabic}

This document ...

\section{Functional Requirements Evaluation}
\subsection{Area of Testing 1: Admin Portal}
\subsubsection{Test-FR-AP-1: Create Form}
\textbf{Type:} Manual \\ \textbf{Initial State:} Administrator is logged into the admin portal with permission to create forms. \\ \textbf{Test Procedure:}
\begin{verbatim}
Input:
- Navigate to the form creation page
- Enter form name = "Event Feedback Form"
- Add field: Text field labelled "Name"
- Add field: Checkbox labelled "Events Attended"
- Add field: Rating field labelled "Overall Experience"
- Click "Create Form"
\end{verbatim}
\textbf{Expected Result:}
\begin{verbatim}
Output:
- System confirms successful form creation
- Newly created form appears in the admin dashboard
- Form fields are stored correctly in the database
\end{verbatim}
\textbf{Result:} Pass \\ \textbf{Discussion:} The administrator successfully created a new form with multiple fields and the form appeared correctly in the dashboard with all fields stored in the system database. \\
\subsubsection{Test-FR-AP-2: View and Analyze Form Data}
\textbf{Type:} Manual \\ \textbf{Initial State:} At least two forms exist with previously collected responses. \\ \textbf{Test Procedure:}
\begin{verbatim}
Input:
- Navigate to the analytics dashboard
- Select multiple forms with collected responses
- Click "View Analytics"
\end{verbatim}
\textbf{Expected Result:}
\begin{verbatim}
Output:
- Aggregated analytics are displayed
- Charts and response summaries appear
- Displayed analytics match stored backend data
\end{verbatim}
\textbf{Result:} Pass \\ \textbf{Discussion:} The analytics dashboard correctly aggregated the responses from multiple forms and displayed charts and summary statistics consistent with the stored backend data. \\
\subsubsection{Test-FR-AP-3: Manage Events via Dashboard}
\textbf{Type:} Manual \\ \textbf{Initial State:} Multiple events exist in the system with various payment statuses. \\ \textbf{Test Procedure:}
\begin{verbatim}
Input:
- Navigate to the admin dashboard
- Select the filter option
- Apply filter "Pending Payment"
\end{verbatim}
\textbf{Expected Result:}
\begin{verbatim}
Output:
- Dashboard displays only events with pending payment status
- Displayed events match stored system records
\end{verbatim}
\textbf{Result:} Pass \\ \textbf{Discussion:} The dashboard filtering functionality worked correctly and displayed only the events matching the selected payment status filter. \\
\subsubsection{Test-FR-AP-4: View and Export Event Analytics}
\textbf{Type:} Manual \\ \textbf{Initial State:} Events exist with collected survey responses. \\ \textbf{Test Procedure:}
\begin{verbatim}
Input:
- Navigate to the event analytics dashboard
- Select an event with stored responses
- Click "Export Analytics"
\end{verbatim}
\textbf{Expected Result:}
\begin{verbatim}
Output:
- Event analytics are displayed in the dashboard
- System generates a downloadable report file
- Exported data matches the analytics displayed
\end{verbatim}
\textbf{Result:} Pass \\ \textbf{Discussion:} The export functionality successfully generated a downloadable analytics report containing the correct event statistics and survey response summaries. \\


\section{Nonfunctional Requirements Evaluation}

\subsection{Usability}
		
\subsection{Performance}

\subsection{etc.}
	
\section{Comparison to Existing Implementation}	

This section will not be appropriate for every project.

\section{Unit Testing}

\section{Changes Due to Testing}

\wss{This section should highlight how feedback from the users and from 
the supervisor (when one exists) shaped the final product.  In particular 
the feedback from the Rev 0 demo to the supervisor (or to potential users) 
should be highlighted.}

\section{Automated Testing}
		
\section{Trace to Requirements}
		
\section{Trace to Modules}		

\section{Code Coverage Metrics}

\bibliographystyle{plainnat}
\bibliography{../../refs/References}

\newpage{}
\section*{Appendix --- Reflection}

The information in this section will be used to evaluate the team members on the
graduate attribute of Reflection.

\input{../Reflection.tex}

\begin{enumerate}
  \item What went well while writing this deliverable? 
  \item What pain points did you experience during this deliverable, and how
    did you resolve them?
  \item Which parts of this document stemmed from speaking to your client(s) or
  a proxy (e.g. your peers)? Which ones were not, and why?
  \item In what ways was the Verification and Validation (VnV) Plan different
  from the activities that were actually conducted for VnV?  If there were
  differences, what changes required the modification in the plan?  Why did
  these changes occur?  Would you be able to anticipate these changes in future
  projects?  If there weren't any differences, how was your team able to clearly
  predict a feasible amount of effort and the right tasks needed to build the
  evidence that demonstrates the required quality?  (It is expected that most
  teams will have had to deviate from their original VnV Plan.)
\end{enumerate}

\end{document}